%%% Изображения %%%
\graphicspath{{images/}{Synopsis/images/}}         % Пути к изображениям

%%% Макет страницы %%%
% symmetric synopsis geomery:
% \geometry{a4paper, top=20mm, bottom=20mm, inner=25mm, outer=10mm, footskip=5mm, nomarginpar}%, showframe
% dissertation geomery:
\geometry{a4paper, top=2cm, bottom=2cm, left=2.5cm, right=1cm, nofoot, nomarginpar} %, heightrounded, showframe
\setlength{\topskip}{0pt}   %размер дополнительного верхнего поля

%%% Интервалы %%%
%% Реализация средствами класса (на основе setspace) ближе к типографской классике.
%% И правит сразу и в таблицах (если со звёздочкой)
%\DoubleSpacing*     % Двойной интервал
% \OnehalfSpacing*    % Полуторный интервал
\SingleSpacing      % Одинарный интервал
%\setSpacing{1.42}   % Полуторный интервал, подобный Ворду (возможно, стоит включать вместе с предыдущей строкой)
% Отделять формулы от текста сверху и снизу одним строковым интервалом
% (устанавливаем величины вокруг display-math в начале документа)
\AtBeginDocument{%
    \setlength{\abovedisplayskip}{1.5\baselineskip}%
    \setlength{\belowdisplayskip}{1.5\baselineskip}%
    \setlength{\abovedisplayshortskip}{1.5\baselineskip}%
    \setlength{\belowdisplayshortskip}{1.5\baselineskip}%
}
% Один интервал сверху и снизу от рисунков и таблиц (ГОСТ)
% \setlength{\intextsep}{1.5\baselineskip}      % над и под float'ами в середине страницы
% \setlength{\textfloatsep}{1.5\baselineskip}   % между текстом и float'ом вверху/внизу страницы

% \setlength{\abovecaptionskip}{\baselineskip}   % space between figure and caption
% \setlength{\belowcaptionskip}{\baselineskip}   % space after caption

\setlength{\parindent}{1.25cm}
\setlength{\parskip}{0pt}

%%% Выравнивание и переносы %%%
%% http://tex.stackexchange.com/questions/241343/what-is-the-meaning-of-fussy-sloppy-emergencystretch-tolerance-hbadness
%% http://www.latex-community.org/forum/viewtopic.php?p=70342#p70342
\tolerance 1414
\hbadness 1414
% \emergencystretch 0.5em % В случае проблем регулировать в первую очередь
\hfuzz 0.3pt
\vfuzz \hfuzz
% % Отключаем автоматические переносы
% % Устанавливаем очень высокий штраф за перенос и за явные дефисы
% \hyphenpenalty=10000
% \exhyphenpenalty=10000
%\raggedbottom
\sloppy                 % Избавляемся от переполнений
\clubpenalty=10000      % Запрещаем разрыв страницы после первой строки абзаца
\widowpenalty=10000     % Запрещаем разрыв страницы после последней строки абзаца

%%% Колонтитулы %%%
\makeevenhead{plain}{}{\thepage}{}
\makeoddhead{plain}{}{\thepage}{}
\makeevenfoot{plain}{}{}{}
\makeoddfoot{plain}{}{}{}
\pagestyle{plain}

%%% Размеры заголовков %%%
\setsecheadstyle{\normalfont\bfseries}
\renewcommand*{\chaptitlefont}{\normalfont\bfseries}

%%% Подписи %%%
\setfloatadjustment{table}{%
    \setlength{\abovecaptionskip}{0pt}   % Отбивка над подписью
    \setlength{\belowcaptionskip}{0pt}   % Отбивка под подписью
}

%%% Отступы у плавающих блоков %%%
% Используем такой же порядок, как и выше (1.5 межстрочного интервала),
% чтобы после таблиц и рисунков был заметный отступ, как после формул.
% \setlength\textfloatsep{1.5\baselineskip}

\newcommand{\theintvl}{2.5} % Интервал между заголовками и текстом в единицах размера шрифта (ГОСТ Р 7.0.11-2011, 5.3.5)
%%% Интервалы между заголовками
\setbeforesecskip{\theintvl\curtextsize}% Заголовки отделяют от текста сверху и снизу тремя интервалами (ГОСТ Р 7.0.11-2011, 5.3.5).
\setaftersecskip{\theintvl\curtextsize}
\setbeforesubsecskip{\theintvl\curtextsize}
\setaftersubsecskip{\theintvl\curtextsize}
\setbeforesubsubsecskip{\theintvl\curtextsize}
\setaftersubsubsecskip{\theintvl\curtextsize}
